\section{Block observations and the free kernels}
\label{sec:free}

The exact integration in Section~\ref{sec:rg} uses a positive decomposition
of the quadratic energy and exponentially localized Gaussian kernels.
We construct these objects for ordinary square blocks and for one inclined
initial block.  The inclined construction is needed for the rotational
comparison in Section~\ref{sec:continuum}.  Its geometry differs only at the
first step; after sufficiently many common ordinary steps, the two free
precisions approach one another at a rate of order $L^{-2}$ per step.
The linear construction and positive decomposition are those of
\Rref{free:section}.  We prove the changes required by the inclined geometry
and retain the stronger history estimate implicit in the proof of
\Rref{app:free-strip}.

Fix a sufficiently large dyadic integer $L$ and put $L_*=5L$.
We use the scales in Equation~\eqref{setup:scale-definitions}, in
particular $\gamma=(\log L)/(2\pi)$, and the history convention of
Section~\ref{setup:scales}.  Constants without an $L$ subscript
in this section are uniform for all sufficiently large $L$ and all
histories.  Constants denoted by $C_L$ may depend on the fixed $L$.
Only finitely many derivative and exponential-moment orders are required
in any application; these orders are fixed before $L$ is chosen.

\subsection{Cells, weighted means, and the spherical observation}

Use the ordinary weights and spherical observation defined in
Section~\ref{setup:scales}.  Write $\chi$ for their fixed even smooth
one-dimensional bump and $\rho=0.14$ for its support radius after
rescaling a block to side one.  The same bump is used in all the
blockings below.  Coordinates in each lattice plane may be translated.

For the two inclined blockings, let
\[
 O_+=\frac15\begin{pmatrix}3&-4\\4&3\end{pmatrix},
 \qquad O_-=O_+^T,
 \qquad c=\left(\frac12,\frac12\right).
\]
For $O\in\{O_+,O_-\}$ and $Y\in\Z^2$, define
\[
 \mathcal B_Y^O
 =\left(c+L_*OY+O[-L_*/2,L_*/2)^2\right)\cap\Z^2.
\]
The coarse site $Y$ is embedded at the center $c+L_*OY$.
Sample the same product bump in the $O$ coordinates of this cell and
normalize it.  If $w_Y(x)$ denotes either an ordinary or an inclined
weight, the averaging operator is
\[
 (Qu)(Y)=\sum_{x\in\mathcal B_Y}w_Y(x)u_x,
 \qquad w_Y(x)\ge0,\qquad \sum_xw_Y(x)=1.
\]
For a cell of side $l\in\{L,L_*\}$, smooth sampling gives
$\max_xw_Y(x)\le C l^{-2}$.  The centered first moment vanishes exactly.
For an ordinary cell this follows from coordinate reflections; for an
inclined cell it follows from the quarter-turn symmetry proved below.

For inclined cells, define $m_Y$, $\widehat m_Y$, and the independent
constituent tags by Equation~\eqref{setup:unit-mean}, using these weights.
The normalized spherical observation is Equation~\eqref{setup:observation}
with this $Q$.  It is again a probability kernel, since its normalizing
integral depends only on the length of the unit vector $\widehat m_Y$.
The scalar free observation uses the same $Q$ and independent Gaussian
noise of variance one.

\begin{lemma}[Geometry and mean error for the inclined cells]
\label{lem:free-inclined-geometry}
For either choice of $O$, the cells $\mathcal B_Y^O$ partition $\Z^2$,
have exactly $L_*^2$ sites, and are connected by nearest-neighbor edges.
There are no fine sites on their boundaries.  Quarter turns about cell
centers preserve the blocking.  Reflection in a fine coordinate axis
interchanges the two inclined types, with the reflected choice of centers.
Every site of a cell can be joined to its central region by a path in the
cell of length $O(L_*)$.  Moreover, for every spin field, every cell, and
every value of its tag,
\begin{equation}
 |m_Y-\widehat m_Y|
 \le \frac{C}{L_*}\sum_{e\subset\mathcal B_Y^O}r_e,
 \qquad r_{\{x,y\}}=|q_x-q_y|.
 \label{eq:orig-3}
\end{equation}
The ordinary cells satisfy the corresponding estimate with $L$ in place
of $L_*$.
\end{lemma}

\begin{proof}
The two block-translation vectors are integral and their determinant is
$L_*^2$.  For example, for $O_+$ they are $(3L,4L)$ and $(-4L,3L)$.
Multiplying a boundary equation by $5$ gives an integral right-hand side
and a half-integral left-hand side, because the center is $(1/2,1/2)$
and each relevant signed sum of $3$ and $4$ is odd.  Thus no boundary
contains a fine site.  A half-open cell is a fundamental region for its
block-translation lattice, and its fine sites give one representative
of each coset.  The number of sites is therefore its index $L_*^2$.
A quarter turn about $c$ preserves $\Z^2$, the rotated square, and its
product weights.  The same holds about every translated center.
Coordinate reflection changes $O_+$ into $O_-$ and transports the
center and weights accordingly.

For connectivity, use the rotated coordinates relative to the cell
center.  If both coordinates have large absolute value, one of the
four fine coordinate steps decreases both absolute values, each by
at least $3/5$.  If only one coordinate is near a cell boundary, choose
a fine step decreasing that coordinate; the other coordinate has room
for its change, whose absolute value is at most $4/5$.  Repeating these
moves reaches a fixed smaller interior rectangle in $O(L_*)$ steps.
The same moves reach a bounded neighborhood of the center; paths of
bounded length within the interior then join the central fine sites.
All moves remain in the cell.  This also supplies the cell paths used
in the block alignment and coercivity estimates.

It remains to prove the quantitative mean estimate.  In fine coordinates,
the support of the weights lies in the square
\[
 A=\{x\in\Z^2: |x_i-c_i|\le (7/5)\rho L_*,\ i=1,2\}
\]
translated to the cell in question.  This square lies strictly within
the inclined cell: its rotated coordinate radius is at most
$(49/25)\rho L_*=0.2744L_*<L_*/2$.
For each ordered pair $x,y\in A$, join $x$ to $y$ by first moving
horizontally and then vertically inside $A$.  In the weighted average
of these paths, a horizontal edge in row $r$ has total load at most
$\sum_{x:x_2=r}w_Y(x)\le C/L_*$.  A vertical edge in column $s$ has
load at most $\sum_{y:y_1=s}w_Y(y)\le C/L_*$.  Hence
\[
 D_Y:=\sum_{x,y}w_Y(x)w_Y(y)|q_x-q_y|
 \le \frac{C}{L_*}\sum_{e\subset A}r_e.
\]
Since the spins have unit length,
\[
 1-|m_Y|^2
 =\frac12\sum_{x,y}w_Y(x)w_Y(y)|q_x-q_y|^2
 \le D_Y.
\]
On the normalized-mean branch,
$|m_Y-\widehat m_Y|=1-|m_Y|\le1-|m_Y|^2$.
On the fallback branch, $|m_Y|<1/2$ and every possible constituent
spin satisfies
\[
 |m_Y-q_{\tau_Y}|\le1+|m_Y|
 \le2(1-|m_Y|^2)\le2D_Y.
\]
This proves Equation~\eqref{eq:orig-3} uniformly in the tag.
The coordinate-path argument on an ordinary square is identical.
\end{proof}


\begin{corollary}[Quadratic mean error under small chords]
\label{cor:free-small-chord-mean}
For either block geometry, suppose that $r_e\le T_j$ on the within-cell
coordinate paths joining weighted sites in the proof of
Lemma~\ref{lem:free-inclined-geometry}; in particular, this holds if all
nearest-neighbor chords in the cell are at most $T_j$.
For $H\ge H_0(L)$ the normalized-mean branch is forced and
\[
 |m_Y-\widehat m_Y|\le C_L T_j^2.
\]
The bound is uniform in the layer $j$ and the cell.
\end{corollary}

\begin{proof}
For either geometry these paths have length at most $Cl$, where
$l\in\{L,L_*\}$.  Thus any two sites of positive weight satisfy
$|q_x-q_y|\le ClT_j$, and the squared pair identity gives
\[
 1-|m_Y|^2
 =\frac12\sum_{x,y}w_Y(x)w_Y(y)|q_x-q_y|^2
 \le C l^2T_j^2\le C_L T_j^2.
\]
Since $T_j\le H^{-49/100}$, taking $H\ge H_0(L)$ makes the right-hand
side less than $3/4$, uniformly in $j$.  Therefore $|m_Y|>1/2$ and
$|m_Y-\widehat m_Y|=1-|m_Y|\le1-|m_Y|^2$, proving the assertion.
\end{proof}

\subsection{Scalar precisions and their dependence on the history}

We first specify the operator spaces and Fourier convention.  On a
lattice $\Lambda$ identified with $\Z^2$ in its own frame, scalar fields
belong to $\ell^2(\Lambda;\R)$ and oriented edge fields to
$\ell^2(\Lambda;\R^2)$.  The forward gradient and its symbol are
\[
 (d_\mu u)(x)=u(x+e_\mu)-u(x),\qquad
 d_\mu(k)=e^{ik_\mu}-1,
 \qquad D(k)=d(k)^*d(k)=4\sum_{\mu=1}^2\sin^2(k_\mu/2).
\]
Adjoints use counting measure.  At a fixed coarse momentum, write
$\alpha$ for the fine alias index.  The one-cell fine norm is
$l^2\sum_\alpha|u_\alpha|^2$, while the coarse norm is $|V|^2$.
Consequently, if $Qu=\sum_\alpha Q_\alpha u_\alpha$, then
\[
 (Q^*V)_\alpha=l^{-2}Q_\alpha^*V.
\]
Conversely, for a coarse-to-fine map with $(TV)_\alpha=T_\alpha V$,
its adjoint is $T^*u=l^2\sum_\alpha T_\alpha^*u_\alpha$.
The adjoints of the finite matrices $Q_\alpha$ and $T_\alpha$ in these
formulas use the ordinary Euclidean inner product.  Stars at complex
momentum mean analytic continuation of the adjoint on real momentum.

A history $h$ is a finite sequence of ordinary observations,
or an inclined initial observation followed by ordinary observations.
The empty history has $K_{\varnothing}=K_0=D$.  If $Q$ is the next
observation from the fine lattice $\Lambda$ to the coarse lattice $\Lambda'$, its output precision, conditional mean, and conditional
covariance are defined by
\begin{equation}
 (K_{h}')^{-1}
 =\Id+QK_{h}^{-1}Q^*,\qquad
 P_{h}=K_{h}^{-1}Q^*K_{h}',\qquad
 C_{h}=(K_{h}+Q^*Q)^{-1}.
 \label{eq:free-operators}
\end{equation}
These formulas at the massless zero mode are understood by continuation
of the Fourier expressions below.  Equivalently, the first formula
computes the covariance of $Qu$ plus the observation noise.  We write
$K_h,P_h,C_h$ when the particular history is understood.  A prime always
refers to the next lattice, not to differentiation.

To make the localization assertions precise, an operator on one lattice
has an exponential kernel bound if, for some $c>0$,
$\sup_x\sum_y e^{c|x-y|}|F(x,y)|<\infty$.
For a map between the fine and coarse lattices, use their physical
positions measured in coarse lattice units.  The analogous supremum
over input sites is its column bound.  Polynomial factors of any fixed
order may be included by reducing $c$.  Fine differences mean
nearest-neighbor differences in the original fine coordinate directions.

\begin{proposition}[Uniform free bounds and history contraction]
\label{prop:free-bounds}
There are $L_0,A_0,c,C>0$ such that, for every dyadic $L\ge L_0$ and
every admissible history, $K_{h}$ is analytic and bounded on a
common complex neighborhood of the momentum torus and
\[
 cD(p)\le K_{h}(p)\le CD(p)\quad(p\in\R^2),
 \qquad K_{h}(p)-D(p)=O(|p|^4).
\]
There is a real, self-adjoint analytic edge lift $B_{h}$ with
\[
 K_{h}=d^*B_{h}d,
 \qquad c\Id\le B_{h}\le C\Id,
 \qquad B_{h}(0)=\Id.
\]
The lifts and their inverses have uniform exponential kernel bounds.
For a step of side $l=L$ or $l=L_*$, the operators in
Equation~\eqref{eq:free-operators} satisfy
\[
 KP=Q^*K',\qquad \Id-QP=K',
 \qquad
 \sup_x\sum_Y e^{c\operatorname{dist}(x,Y)/l}
       |\nabla_f^sP(x,Y)|\le C_s l^{-s}
\]
for every fixed nonnegative integer $s$.  Here the distance is between
the physical fine site and coarse center, in fine units.
The mean $P$ preserves constants and centered affine fields, with the
rotation and scale of the output frame included.  The covariance $C$
has exponential kernel bounds in coarse units with constants $C_L$.

Suppose two histories share their last $r\ge0$ ordinary observation
steps, using the same ordinary $Q$ in these steps and identifying their
output coordinates.  Then on a common fixed smaller complex strip,
\begin{equation}
 \norm{K_{h_1}-K_{h_2}}_{\rm an}
 \le C\left(\frac{A_0}{L^2}\right)^r.
 \label{eq:free-history}
\end{equation}
Here $\norm{\cdot}_{\rm an}$ is the supremum norm on that strip.
The difference vanishes through degree three at the origin, and the
same estimate holds for any prescribed finite set of analytic derivative
and exponential kernel norms, after decreasing the strip or exponential
weight.  The corresponding differences of lifts and their inverses
obey this bound.  When the two histories are followed by a common next
ordinary observation, their $P$ and $C$ differences obey the same bound
with a permitted prefactor $C_L$.
\end{proposition}

\begin{proof}
We give the analytic estimates with the geometry visible.  This also
identifies the hypotheses of \Rref{free:bounds} that are retained.
In an inclined step of side $l=L_*$, the fine aliases above coarse
momentum $p$ are
\[
 k_\alpha=O(p+2\pi\alpha)/l,
 \qquad \alpha\in\Z^2/(lO^T\Z^2).
\]
Choose representatives in a centered rotated square.  For an ordinary
step use $O=\Id$ and $l=L$.  Write
$b_Q(k)=\sum_xw_0(x)e^{ik\cdot x}$, with integer origins when aliases
are glued.  Centered origins only change this symbol by a translation
phase.  On
\[
 \Omega_\delta=\{p\in\C^2:
       |\operatorname{Re}p_\mu|<\pi+\delta,       |\operatorname{Im}p_\mu|<\delta,\ \mu=1,2\},
\]
discrete summation by parts in the fine coordinate directions gives,
for any fixed $J$ and multi-index $\beta$,
\begin{equation}
 |\partial_p^\beta b_Q(k_\alpha)|
 \le C_{J,\beta}(1+|\alpha|)^{-J}.
 \label{eq:free-bump-decay}
\end{equation}
Indeed, the sum of absolute $J$th differences of the sampled normalized
weights is $O(l^{-J})$, while at least one fine difference multiplier
on a nonprincipal alias has size $c(1+|\alpha|)/l$.  The weights vanish
near the cell boundary, so the summation produces no boundary term.
Complex tilting and $p$ derivatives insert bounded rescaled coordinates.
This is the proof of \Rref{app:analytic-1}, without any restriction to
an axis-aligned support.

The centered principal bump is uniformly nonzero on the real principal
cube.  For ordinary blocks use its product structure and the small
support.  For inclined sampling, its Riemann sum converges uniformly on
that cube to the product Fourier integral in the rotated coordinates;
that integral is nonzero by the same support bound.  Increasing $L_0$
and decreasing the complex strip preserve the lower bound.
The exact first moment is zero, so the principal product has the
quadratic Taylor bound used in \Rref{app:free-strip}.
Representatives relabel on overlaps of momentum cubes; the centered
representatives there have comparable alias weights.  We use these
overlaps also at the faces of the rotated alias box.

For a compound history let $m$ be the product of its side factors
and let $O_h$ map the output coordinate frame to the original fine
frame.  For the coarse site at the origin, let $v_h(x)$ be the product
of observation weights along the unique ancestry of its fine
descendant $x$.  These weights sum to one.  Write
$c_h=\sum_xv_h(x)x$ for their center in the original fine coordinates.
The transform in centered output coordinates is
\[
 W_h(\xi)=\sum_xv_h(x)
   \exp\left\{i\xi\cdot O_h^T(x-c_h)/m\right\},
 \qquad D_{m,O_h}(\xi)=m^2D(O_h\xi/m).
\]
Using integer origins instead multiplies $W_h$ by a translation phase,
which cancels in $W_hW_h^*$.  This is the origin convention used when
gluing aliases.
Separating the principal alias in the covariance gives precisely the
formula of \Rref{free:explicit}:
\begin{equation}
 K_{h}(p)=
 \frac{D_{m,O_{h}}(p)}
 {W_{h}(p)W_{h}(p)^*
       +D_{m,O_{h}}(p)R_{h}(p)}.
 \label{eq:free-explicit}
\end{equation}
Here $R_{h}$ is the sum over nonprincipal aliases of
$W_{h}W_{h}^*/D_{m,O_{h}}$, together with
the independent noise variances.  The latter form a sum of products
of squared-weight sums, starting with the last observation.  This
order matters when the first observation is inclined.
For the empty history, take $W=1$ and $R=0$.

For every nonprincipal centered representative, on a fixed sufficiently
small complex domain,
\begin{equation}
 |D_{m,O_{h}}(p+2\pi\alpha)|
 \ge c(1+|\alpha|)^2.
 \label{eq:free-alias-denominator}
\end{equation}
To verify this, use $D(z)=4\sum_\mu\sin^2(z_\mu/2)$ in the fine
coordinates.  On the centered fine box its real part is bounded below
by a fixed multiple of the squared real argument minus a fixed
multiple of the squared imaginary argument.  A nonprincipal alias
has real distance at least $c(1+|\alpha|)$ from the principal origin.
The same conclusion holds for representatives slightly across a face,
by periodicity.  Thus the proof of \Rref{app:analytic-2} applies with
the indicated rotation of the argument.

The remaining issue is a bound independent of history length.  In a
sequence of $i$ ordinary observations, each descendant has a unique
ancestral position and the squared weights sum to at most
$(A/L^2)^i$, for a fixed $A$.  Descendant positions divided by their
total side remain in a bounded set.  Parseval at opposite imaginary
tilts, followed by Cauchy--Schwarz, therefore gives the estimate of
\Rref{app:analytic-3}:
\[
 \sum_{\alpha\bmod L^i}
   |W_i(p+2\pi\alpha)W_i(p+2\pi\alpha)^*|
 \le C A^i.
\]
For a ball of radius $C_1L^i$, the same bound holds with a changed
fixed constant, since that ball meets only a bounded number of alias
boxes.  Earlier factors have absolute product bounded by a fixed
constant: their real arguments cost at most one, and their complex
tilts have a summable geometric bound.  Split the aliases into shells
between successive radii $c_1L^i$.  By
Equation~\eqref{eq:free-alias-denominator}, their contributions have
successive ratio at most $A/L^2$, after changing $A$ once.  If the
history begins with an inclined observation, use Parseval also on the
final shell ending at its full rotated box.  Its side factor $5L$
changes only fixed constants.  This proves the shell bound of
\Rref{app:analytic-4}, uniformly for both kinds of history.
The noise sum is bounded by the same geometric series.

Consequently $R_{h}$ and its required derivatives are
uniformly bounded on a fixed strip.  On the real principal cube,
$W_{h}W_{h}^*$ is uniformly bounded below: its first
factor has the lower bound already proved, and the departures of
subsequent factors from one have a summable quadratic bound.
The denominator in Equation~\eqref{eq:free-explicit} is therefore
positive and uniformly bounded below on the real cube.  Its uniformly
bounded derivatives keep it nonzero on a smaller fixed complex strip.
This proves analyticity and ellipticity.  At the origin its denominator
is $1+O(|p|^2)$, while its numerator has quadratic term $|p|^2$.
Reality and inversion symmetry remove odd terms.  Every precision
therefore has the same jet through degree three.

We now compare two histories with $r$ common final ordinary steps.
For a sufficiently small fixed $c_2$, aliases with
$|\alpha|<c_2L^r$ can be identified in both histories.  Put
$\xi=p+2\pi\alpha$ and denote the common product of the last $r$
factors by $W_r$.  Taylor expansion of all factors preceding these
steps, with the linear terms canceled in each $bb^*$, gives
\[
 |W_{h_1}W_{h_1}^*
       -W_{h_2}W_{h_2}^*|
 \le C|\xi|^2L^{-2r}|W_rW_r^*|.
\]
No division by a bump factor is used.  On the common nonprincipal
aliases, the inverse denominators differ by at most $CL^{-2r}$.
For the rotated denominator this follows from the same Taylor
estimate, since its leading quadratic form is still $|\xi|^2$.
Parseval now bounds the change of the common nonprincipal sum by
$CA^rL^{-2r}$.  The shells outside this common region cost at most
$CA^{r+1}L^{-2r}$, and the part of the noise sum preceding the common
steps has the same bound.  On the principal term apply Taylor's
estimate directly to the numerator and product in
Equation~\eqref{eq:free-explicit}.  Its denominator remains nonzero
on the common strip, giving Equation~\eqref{eq:free-history} after
increasing $A_0$.  This argument also covers $r=0$, when only a uniform
bound is asserted.  Cauchy estimates on nested fixed strips give all
prescribed derivatives, and Fourier inversion gives the exponential
kernel estimates and their polynomial moments.

For completeness, the edge lift is obtained by applying the bounded
coordinate divisions of \Rref{app:free-strip} to $K-D$, which vanishes
to order four.  They give a Hermitian lift
$B^{\rm raw}=\Id+O(|p|^2)$ with bounded analytic coefficients.
If $\mathsf c=(-d_2,d_1)$, then $\mathsf c d=0$, so addition of a
fixed sufficiently large multiple of $\mathsf c^*\mathsf c$ leaves
$d^*Bd$ unchanged.  The longitudinal form is bounded below by $K/D$;
a two-by-two Schur complement makes the resulting $B$ uniformly
positive away from zero, and $B^{\rm raw}$ is already positive near
zero.  All divisions are bounded linear operations on fixed nested
strips.  The common added term cancels in differences, and the inverse
identity preserves Equation~\eqref{eq:free-history}.  At empty history
one may keep $B=\Id$.

In aliases the conditional mean is
\[
 P_\alpha=l^{-2}K(k_\alpha)^{-1}b_Q(k_\alpha)^*K'(p).
\]
On nonprincipal aliases, Equations~\eqref{eq:free-bump-decay}
and~\eqref{eq:free-alias-denominator} give arbitrary fixed inverse-power
alias decay.  The principal pole is removable by the covariance
identity; its zero and first jets prove preservation of constants and
centered affine fields.  In inclined coordinates an affine function
is evaluated at $x=c+lOY$, which is the rotation and rescaling in
\Rref{supp:fixed-norm-affine-assumption}.  Each fine difference
contributes at most $C(1+|\alpha|)/l$.  Taking sufficiently many bump
differences and Fourier inverting on a smaller strip proves the
claimed bounds for every fixed order $s$, including differences
across cell boundaries.  Finally, the nonprincipal block of
$K+Q^*Q$ is inverted first; its rank-one update and principal Schur
complement have the positive denominators used in \Rref{free:bounds}.
This proves the covariance bound, allowing constants depending on $L$.
All these operations preserve the history discrepancy for a common
next observation.  For occurrences of $K^{-1}$ on nonprincipal aliases,
use the inverse difference identity and the same denominator bounds.
\end{proof}

\subsection{The positive kinetic identity and symmetry of its construction}

The scalar estimates do not by themselves provide the positive local
energy used in the nonlinear integration.  We next lift them to edge
operators.  Write
$\mathcal X_\Lambda=\ell^2(\Lambda;\R^2)\oplus
\ell^2(\Lambda;\R^2)$ for a pair of edge fields.
The following identity applies componentwise also to ambient
$\R^3$-valued fields; no constraint on their values is needed.

\begin{proposition}[Positive completion for both block geometries]
\label{prop:free-positive}
A common $c_0>0$ and real analytic edge operators $\ell_{h}$
can be chosen so that
\[
 \mathsf A_{h}
   =\binom{\sqrt{c_0}\Id}{\ell_{h}},
 \qquad X_{h}=\mathsf A_{h}d,
 \qquad X_{h}^*X_{h}=K_{h}.
\]
There are real maps
\[
 \mathcal M:\mathcal X_{\rm coarse}
        \longrightarrow\mathcal X_{\rm fine}\oplus\ell^2(\Lambda';\R),
 \qquad
 \mathcal N:\mathcal X_{\rm coarse}
        \longrightarrow\mathcal X_{\rm coarse}\oplus\ell^2(\Lambda';\R)
\]
with
\begin{align}
 \mathcal M^*\mathcal M+\mathcal N^*\mathcal N&=\Id,
 &\mathcal N X'&=0,\label{eq:free-positive-isometry}\\
 \mathcal M^*(Xu,V-Qu)&=X'V,
 &\mathcal M X'&=(XP,\Id-QP).
 \label{eq:free-positive-ambient}
\end{align}
Their kernels have exponential moments.  For a block side
$l\in\{L,L_*\}$, the fine-output part of $\mathcal M$ has row bound
$C/l$, column bound $Cl$, and first and second fine-difference row
bounds $C/l^2$ and $C/l^3$.  The coarse slots have bounded row and
column norms.  For a common next ordinary step, every history difference
has the factor $(A_0/L^2)^r$ of Equation~\eqref{eq:free-history},
with prefactor $C_L$ allowed.  At empty history one may take
$\ell_{\varnothing}=\sqrt{1-c_0}\Id$.

The construction commutes with simultaneous square symmetries of its
geometry and its input data.  An individual inclined history has C4
covariance; reflection transports it to the reflected history.  Ordinary
histories, including empty history, have the full square symmetries.
\end{proposition}

\begin{proof}
The construction in \Rref{lem:free-stack} begins with an operator $T$
from coarse edges to fine edges satisfying
\begin{equation}
 T^*d=d'Q,\qquad TB'd'=BdP.
 \label{eq:free-edge-identities}
\end{equation}
We verify both analytic divisions in that construction for the inclined
step.  For each block-translation vector $lOe_\mu\in\Z^2$, choose a
fine coordinate path from $0$ to that vector.  Its edge Fourier row
$\omega_\mu(k)$ has length $O(l)$ and satisfies
\[
 \omega_\mu(k)d(k)=e^{ik\cdot lOe_\mu}-1=d'_\mu(p).
\]
Collect these rows in $\omega$ and set
$T_{0,\alpha}=l^{-2}\omega(k_\alpha)^*b_Q(k_\alpha)^*$.
The fine adjoint normalization gives $T_0^*d=d'Q$.
For an ordinary step this is the coordinate-path formula in
\Rref{lem:free-stack}.

The remainder
$R_\alpha=B_\alpha d_\alpha P_\alpha-T_{0,\alpha}B'd'$
is divergence free, by $KP=Q^*K'$.
The principal values of the path rows are the translation vectors,
while $d_0P_0$ has linear term $iOp/l$.  These linear terms cancel
in $R_0$, so it vanishes through first order.  On every nonprincipal
alias $R_\alpha(0)=0$.  Its norm is $O(l^{-1})$ with arbitrary
prescribed inverse-power decay in the alias index.
The first division of \Rref{app:alias-division} therefore gives
\[
 R_\alpha=\mathsf c(k_\alpha)^*\gamma_\alpha,
 \qquad \gamma_\alpha(0)=0.
\]
At the principal origin use local coordinates along the fine axes;
the rotation changes constants by a fixed factor.  On a nonprincipal
overlap divide by a component of $d(k_\alpha)$ of size at least
$c(1+|\alpha|)/l$.  The resulting scalar quotients agree on overlaps
by the divergence-free identity.  Thus $\gamma$ has the same arbitrary
fixed alias decay, without the factor $l^{-1}$.

The second division must produce a row $H$ with $Hd'=\gamma$ while
preserving this decay.  In the inclined case the true momentum-period
lattice, in the $p$ axes and with the factor $2\pi$ omitted, is
\[
 lO^T\Z^2
 =L\begin{pmatrix}3&\phantom{-}4\\-4&3\end{pmatrix}\Z^2
\]
for $O_+$, and its reflected version for $O_-$.
The rectangular lattice $(25L)\Z^2$ is a sublattice of index $25$.
Work first on this rectangular covering torus.  Its alias decay weight
is the sum of the $25$ translated weights centered at the copies of
the principal alias.  Apply the sinc-power interpolation operator
$\Pi$ of \Rref{app:analytic-5}, with $25L$ in place of the rectangular
side.  For each translated weight its cyclic convolution bound is
unchanged.  The two rows
\[
 H=\left(\frac{\gamma-\Pi\gamma}{d'_1},
          \frac{\Pi\gamma}{d'_2}\right)
\]
are analytic: the first numerator vanishes on each $d'_1=0$ plane;
the second vanishes on each $d'_2=0$ plane because $\gamma$ vanishes
at every sampling intersection.  Cauchy estimates in the $p$ variables
give bounded removable quotients, with no extra power of $l$.
Average over the finite deck group to recover the true period lattice.
Integer fine-site phases make this descent agree with alias gluing.
We have therefore obtained the second division with constants uniform
in $L$.  Setting
\[
 T_\alpha=T_{0,\alpha}
       +\mathsf c(k_\alpha)^*H_\alpha(B')^{-1}
\]
proves Equation~\eqref{eq:free-edge-identities}, with alias bound
$C_Jl^{-1}(1+|\alpha|)^{-J}$.
Fourier inversion gives its row, column, and fine-difference bounds.
The divisions and the inverse identity preserve history differences.

We describe the remaining positive completion to specify the role of
the normalization.  On real momentum put
\[
 U=(B')^{-1}-T^*B^{-1}T-d'd'^*.
\]
The edge identities and $\Id-QP=K'$ give
\[
 UB'd'=d'-T^*dP-d'K'=d'(\Id-QP-K')=0.
\]
Coarse coordinate-plane division on both sides factors $B'UB'$
through $\mathsf c'^*\mathsf c'$ with a bounded analytic scalar
coefficient.  Also $\mathsf cB^{-1}TB'd'=0$.
Before the last division its aliases have bound
$Cl^{-2}(1+|\alpha|)^{-J}$.  Squaring and summing with the fine norm
$l^2\sum_\alpha$ gives exactly the estimates used in
\Rref{lem:free-stack}:
\[
 |z^*Uz|\le C|\mathsf c'(B')^{-1}z|^2,
 \qquad
 \norm{\mathsf cB^{-1}Tz}^2
 \le Cl^{-2}|\mathsf c'(B')^{-1}z|^2.
\]
In particular the inclined step has the correct normalization
$l^2=L_*^2$, rather than the area of an enclosing axis square.

For one fixed sufficiently large $\lambda$, replace the inverse lift by
\[
 \widetilde B^{-1}
 =B^{-1}+\lambda B^{-1}\mathsf c^*\mathsf cB^{-1}
\]
and make the same replacement at the next scale.
It leaves $\widetilde Bd=Bd$ and hence the precision unchanged.
Define the modified defect by
\[
 U_*=(\widetilde B')^{-1}
          -T^*\widetilde B^{-1}T-d'd'^*.
\]
Expanding the two modified inverse lifts gives
\[
 U_*=U+\lambda\bigl[
 (B')^{-1}\mathsf c'^*\mathsf c'(B')^{-1}
 -T^*B^{-1}\mathsf c^*\mathsf cB^{-1}T\bigr].
\]
The preceding two bounds therefore show that $U_*$ is bounded below by
\[
 [\lambda(1-C/l^2)-C]
 (B')^{-1}\mathsf c'^*\mathsf c'(B')^{-1}.
\]
Choose $\lambda$ first and then $L$ sufficiently large, so this
coefficient is at least one.  The defect has a factorization
$\widetilde B'U_*\widetilde B'=\mathsf c'^*B_2\mathsf c'$ with
$B_2$ a uniformly positive analytic scalar.  Positivity through zero
follows from
\[
 \mathsf c'(B')^{-1}\widetilde B'
 =\bigl[1+\lambda\mathsf c'(B')^{-1}\mathsf c'^*\bigr]^{-1}
   \mathsf c'.
\]
Choose $c_0$ below the common lower bound for $\widetilde B$ and put
$\ell=(\widetilde B-c_0\Id)^{1/2}$.
For $S'=\mathsf A'(\widetilde B')^{-1}(\mathsf A')^*$ define
\begin{align*}
 \mathcal Mz
 &=\left(\mathsf A\widetilde B^{-1}T(\mathsf A')^*z,
                       d'^*(\mathsf A')^*z\right),\\
 \mathcal Nz
 &=\left((\Id-S')z,
       B_2^{1/2}\mathsf c'(\widetilde B')^{-1}(\mathsf A')^*z\right).
\end{align*}
The identities $\mathsf A^*\mathsf A=\widetilde B$ and
$(\mathsf A')^*\mathsf A'=\widetilde B'$ make $S'$ an orthogonal
projection.  The definition and factorization of $U_*$ then give
\[
 \mathcal M^*\mathcal M
   =S'-\mathsf A'U_*(\mathsf A')^*,\qquad
 \mathcal N^*\mathcal N
   =\Id-S'+\mathsf A'U_*(\mathsf A')^*.
\]
Their sum proves the isometry in
Equation~\eqref{eq:free-positive-isometry}.  Substitution of the two
edge identities proves Equation~\eqref{eq:free-positive-ambient},
and $\mathsf c'd'=0$ proves the null identity.
All square roots and inverses are uniformly analytic on a smaller
strip.  Their exponential kernel bounds preserve those of $T$.
At empty history retain $B=\widetilde B=\Id$ on the fine side;
this removes a nonnegative subtraction from the defect and improves
its positivity.  It gives the asserted value of $\ell_{\varnothing}$.

Finally the choices of lift and intertwiner must be the same rules
for all data, including data with different individual reflection
symmetries.  Given a particular construction $F$ and input data
$\mathcal D$, replace it by the finite average
\[
 \frac1{|G|}\sum_{R\in G}R^{-1}F(R\mathcal D),
\]
where $G$ is the square symmetry group and the actions include the
induced edge and plaquette permutations, signs, and integer
translations restoring the chosen origins.  First make this average
for the lift; then make it for $T$ with that lift fixed.
Each defining identity is linear in the output being averaged, and
positivity and lower bounds of lifts survive the average.
For inclined data include both reflected types in the input family.
The subsequent factorization and positive square roots commute with
these actions.  Thus the rules are covariant even when an individual
precision has only C4 symmetry.  The averaging never replaces that
precision by a mirror average.  Being fixed finite averages, these
operations also preserve all discrepancy estimates.  This common
choice is used in the nonlinear comparisons
\eqref{eq:orig-5} and~\eqref{eq:orig-29}.
\end{proof}

\subsection{Finite periods and unfolded symmetry}

All kernels above are chosen once on the infinite lattice and then
periodized.  The permitted finite tori are quotients by square,
rectangular, or oblique sublattices of translations, with the bounded
shape ratios used below.  Their period lattices must be divisible by
every applied blocking: after each block map they become translation
periods of the next lattice.  We require only that the shortest period,
measured in layer-$0$ sites, exceed a fixed $m_{\min}(L,H)$.

Periodization folds the absolutely summable bulk kernels and their
identities.  For expressions carrying connecting paths, retain the
lifted path before folding, as in \Rref{free:section}; its length
records a contribution that crosses a period.  Locality tests at any
layer use the shortest period in that layer.  Enlarging
$m_{\min}(L,H)$ enforces simultaneously all scale-neighborhood
conditions following \Rref{rg:scales}, including those for the inclined
initial step.  Calculations on short symmetry orbits use C4, or the
full square group when available, on the unfolded bulk labels.
They therefore remain valid on asymmetric finite tori.  Terms whose
records cross a period keep their separate long-path bounds; no
symmetry of such terms, or of the finite period lattice itself, is
assumed.
